\documentclass[11pt,a4paper]{article}
\usepackage[T1]{fontenc}
\usepackage{lmodern,amsmath,amssymb,amsthm,mathtools,booktabs,microtype}
\usepackage[margin=26mm]{geometry}
\usepackage[colorlinks=true,linkcolor=blue,citecolor=blue,urlcolor=blue]{hyperref}
\setlength{\emergencystretch}{2em}
\newtheorem{theorem}{Theorem}[section]
\newtheorem{lemma}[theorem]{Lemma}
\newtheorem{proposition}[theorem]{Proposition}
\newtheorem{corollary}[theorem]{Corollary}
\newcommand{\Q}{\mathbb Q}\newcommand{\Z}{\mathbb Z}
\DeclareMathOperator{\Tr}{Tr}\DeclareMathOperator{\Nm}{Nm}
\title{Quadratic reference covers and weighted numerator charts}
\author{}\date{16 November 2026}
\begin{document}\maketitle
\begin{abstract}
We give two polynomially isomorphic quadratic presentations of a
square-root cover of an elliptic point and its translate. One uses
three field coordinates; the other uses four and exposes both
reciprocal abscissas and two scaled ordinates as linear coordinates.
The inverse maps have only constant denominators and include the
four points above the identity and the translating point. A homogeneous
input scaling preserves quadratic equations. Three field coordinates
then describe a weighted numerator point, and one additional field
coordinate records its physical ratio by a quadratic equation.
Separately, a Jacobi addition chart recovers the original difference
trace with one additional rational coordinate; its zero-scale fiber
is classified exactly. The algebraic interfaces and their coordinate
costs are explicit. They do not yet give a complete defining quartic
with fewer variables.
\end{abstract}

\section{The reference cover}
Let $k$ have characteristic zero. Fix $e,c\in k$ such that
\[
 a=e^2,\qquad b=ac,\qquad D=a(a-c^2),\qquad ecD\ne0.
\]
Thus $P=(a,b)$ lies on $E:y^2=x^3-Dx$. Put $T_0=(0,0)$.
Consider the curve $\mathcal C_3$ with coordinates $r,r',z$ and
two quadratic equations
\begin{equation}\label{eq:three}
 z=rr',\qquad a(r^2+r'^2)+2cez+Dz^2=1.
\end{equation}
The involution $(r,r',z)\mapsto(-r,-r',z)$ has no fixed point.

\begin{theorem}\label{thm:cover}
There is a degree-two etale map
\[
 \pi:\mathcal C_3\longrightarrow E\setminus\{T_0,P+T_0\}.
\]
On $rr'\ne0$ its image $H$ satisfies
\[
 f(H)=r^2,\qquad f(H-P)=r'^2,
 \qquad f(Q)=1/x(Q),
\]
and, with $S_H=y(H)/x(H)$ and $T_H=y(H-P)/x(H-P)$,
\begin{equation}\label{eq:scaledordinates}
 rS_H=e r'+cr+(D/e)rz,\qquad
 r'T_H=-er-cr'-(D/e)r'z.
\end{equation}
The two points $(0,\pm1/e,0)$ map to $O$, and
$(\pm1/e,0,0)$ map to $P$. The map covers every other allowed
reference point exactly when its reciprocal abscissa is a square
in $k$. Its two lifts differ by the displayed involution.
\end{theorem}
\begin{proof}
Use the complete compact reference chart of~\cite{Compact}, obtained
from the translated-point chart of~\cite{Three}. Write
$\delta=D/a$ and set
\[
 \ell=ez,\quad J=\ell^2,\quad
 \alpha=1-ar^2-c\ell,\quad \gamma=\alpha-\delta J-1.
\]
Equations~\eqref{eq:three} give
\[
 \alpha\gamma+c\ell(1-\delta J)+2\delta J=0,
\]
which is the compact reference equation. Its reciprocal functions are
\[
 \frac{1-\alpha-c\ell}{a}=r^2,\qquad
 \frac{\alpha-\delta J-c\ell}{a}=r'^2.
\]
On $\ell\ne0$ its ratios are $S_H=\alpha/\ell$ and
$T_H=\gamma/\ell$. Substitution proves~\eqref{eq:scaledordinates}.
The compact chart is isomorphic to the stated two-point complement,
including $H=O,P$, by~\cite[Theorem 1.1]{Compact}.

Conversely, when $f(H)\ne0$, choose $r$ with $r^2=f(H)$ and put
$r'=\ell/(er)$. The compact equations give $r'^2=f(H-P)$ and
both equations~\eqref{eq:three}. This includes $H=P$. At $H=O$
the two lifts are $(0,\pm1/e,0)$. Near either lift, the derivative
of the second equation with respect to $r'$ is $2ar'\ne0$,
and $\ell=err'$ has a simple zero in the local parameter $r$.
Thus the two lifts are unramified there. Away from $H=O$, the
adjoined nonzero square root has derivative $2r\ne0$ and is etale.
The coordinate ring is generated over the compact reference ring by
$r,r'$, both satisfying monic square equations, so the map is finite.
The local descriptions prove degree two and etaleness everywhere.
No point with $r=r'=0$ satisfies~\eqref{eq:three}, so the two
lifts never coincide.
\end{proof}

The functions in~\eqref{eq:scaledordinates}, denoted below by $s,t$,
satisfy
\begin{equation}\label{eq:quarticordinates}
                 s^2=1-Dr^4,\qquad t^2=1-Dr'^4.
\end{equation}
Thus the point recovery is explicitly
\[
 H=(r^{-2},s/r^3),\qquad H-P=(r'^{-2},t/r'^3)
\]
where the denominators are nonzero, with the stated regular extensions.
This includes tangent configurations and the allowed two-torsion
strata over general fields; no variable ordinate was divided by.

\section{Four coordinates with linear outputs}
Let $\mathcal C_4$ have coordinates $R,Q,s,t$ and equations
\begin{equation}\label{eq:four}
 \begin{split}
 s^2+DR^2&=1,\\
 t^2+DQ^2&=1,\\
 ce\,st&=-a+D(R+Q)+aDRQ.
 \end{split}
\end{equation}

\begin{theorem}\label{thm:isomorphism}
The maps
\begin{equation}\label{eq:forward}
 \begin{split}
 R&=r^2,\qquad Q=r'^2,\\
 s&=er'+cr+(D/e)rz,\\
 t&=-er-cr'-(D/e)r'z
 \end{split}
\end{equation}
and
\begin{equation}\label{eq:inverse}
 \begin{split}
 r&=-\tfrac12\bigl(t/e+eRt+cQs\bigr),\\
 r'&=\tfrac12\bigl(s/e+cRt+eQs\bigr),\\
 z&=\frac{1-a(R+Q)-DRQ}{2ce}
 \end{split}
\end{equation}
are inverse polynomial isomorphisms between $\mathcal C_3$ and
$\mathcal C_4$. In particular every field solution of~\eqref{eq:four}
has $R=r^2$ and $Q=r'^2$ with the roots given by~\eqref{eq:inverse}.
There is no additional nonzero condition on $R,Q,s,t$.
\end{theorem}
\begin{proof}
Substituting~\eqref{eq:forward} into~\eqref{eq:four} and using
\eqref{eq:three} gives the three displayed equations. For example,
\[
 st=-2az-ce(R+Q),\qquad
 2cez=1-a(R+Q)-DRQ
\]
give the third equation. The first two are~\eqref{eq:quarticordinates}.
Reduction by the same equations gives
\[
 r=-\tfrac12(t/e+eRt+cQs),\qquad
 r'=\tfrac12(s/e+cRt+eQs),
\]
so the composition on $\mathcal C_3$ is the identity.

Conversely, substitution of~\eqref{eq:inverse} into
\eqref{eq:three} gives zero modulo the three equations
\eqref{eq:four}. The same calculation gives
$r^2=R$, $r'^2=Q$ and recovers $s,t$ in~\eqref{eq:forward}.
These are polynomial identities with constant denominators, not
identities obtained by dividing by a coordinate. The exact certificates
in Section~\ref{sec:verification} verify the reductions in the original
defining ideals; their parameter denominators have support only on
$e,c,e-c,e+c$, all nonzero under the stated hypotheses.
\end{proof}

The boundary fibers are exactly
\begin{align*}
 R=0:&\qquad Q=1/a,\quad (s,t)=(\pm1,\mp c/e),\\
 Q=0:&\qquad R=1/a,\quad (s,t)=(\pm c/e,\mp1).
\end{align*}
For example setting $R=0$ in~\eqref{eq:four} and eliminating $s,t$
gives $D(aQ-1)^2=0$, after which the signs follow. The other fiber
is symmetric. They are the four points in Theorem~\ref{thm:cover}.

\section{The image over the project quadratic field}
Now put
\[
 K=\Q(\rho),\quad\rho^2=2,\quad d=3-\rho,\quad D=d^2,
 \quad R_0=(1+\rho,2\rho),\quad P=2R_0.
\]
The constants above may be taken as
\[
 e=-1+\frac74\rho,\quad
 a=\frac{57}{8}-\frac72\rho,\quad
 c=-\frac{87}{41}+\frac{129}{164}\rho.
\]
They satisfy $e^2=a$ and $D=a(a-c^2)$. Retain the source group
premise $2E(K)=\Z P$ with $P$ non-torsion, equivalently
$E(K)=\Z R_0+E[2](K)$, as in~\cite{Quartic}.

\begin{proposition}\label{prop:image}
Under this group premise the image of $\mathcal C_3(K)$, or
$\mathcal C_4(K)$, is exactly $2E(K)$. Every such point has two
lifts. A finite point with nonzero abscissa lies in $2E(K)$ exactly
when its abscissa is a square in $K$.
\end{proposition}
\begin{proof}
Use the elementary Kummer homomorphism
\[
 \alpha_E:E(K)\longrightarrow K^\times/K^{\times2},\qquad
 \alpha_E(O)=1,\quad\alpha_E(T_0)=-D,
 \quad\alpha_E(Q)=x(Q)\quad(Q\ne O,T_0).
\]
The chord rule proves multiplicativity: for a nonvertical line
$y=mx+b_0$, the product of the three intersection abscissas is
$b_0^2$. The identity
$x(Q+T_0)=-D/x(Q)$ gives the exceptional torsion rule, and the
duplication formula shows that all doubles lie in the kernel.

The classes of $R_0,T_0,(d,0)$ are respectively
$[1+\rho],[-1],[d]$. They are independent. Indeed the norm of
$(1+\rho)^\epsilon d^\eta(-1)^\zeta$ is
$(-1)^\epsilon7^\eta$. For a square in $K$ this must be a
nonnegative rational square, forcing $\epsilon=\eta=0$.
Formal reality then forces $\zeta=0$. These three classes account
for the eight elements of $E(K)/2E(K)$ under the retained group
premise. Hence $\ker\alpha_E=2E(K)$.

Theorem~\ref{thm:cover} now gives the image statement, including
$O,P$. Neither omitted point $T_0,P+T_0$ belongs to $2E(K)$.
\end{proof}

The norms $\Nm(d)=7$ and $\Nm(1/d)=1/7$ show in particular that
neither $1/d$ nor $-1/d$ is a square in $K$.
Consequently the $K$-points of these square-root covers have no
nonzero-abscissa two-torsion image.

\section{Homogeneous input scaling}
Let $\beta\in\Q$ and retain the rational square coordinates
$b_2=\beta^2$, $b_4=b_2^2$. Use four field coordinates
$R,Q,S,T_H$ satisfying
\begin{equation}\label{eq:homogeneous}
 \begin{split}
 S^2+DR^2&=b_4,\\
 T_H^2+DQ^2&=b_4,\\
 ce\,S T_H&=-ab_4+Db_2(R+Q)+aDRQ.
 \end{split}
\end{equation}
All three equations are quadratic in the displayed coordinates and
the shared power coordinates $b_2,b_4$.

\begin{proposition}\label{prop:scaling}
For $\beta\ne0$, dividing the four coordinates in
\eqref{eq:homogeneous} by $b_2$ gives exactly~\eqref{eq:four}.
If the corresponding reference point is $H$, then
\[
 R=\beta^2 f(H),\quad Q=\beta^2 f(H-P),\quad
 S=\beta^2 rS_H,\quad T_H=\beta^2 r'T_{H,\mathrm{phys}},
\]
where $r^2=f(H)$, $r'^2=f(H-P)$, and the physical ratios are
those of Theorem~\ref{thm:cover}. At $\beta=0$, the only
$K$-solution of~\eqref{eq:homogeneous} is $R=Q=S=T_H=0$.
\end{proposition}
\begin{proof}
Each row is $\beta^4$ times its unscaled row. For the last
assertion, $D=d^2$ is positive at both real embeddings. Thus
$S^2+(dR)^2=0$ forces $S=R=0$, and the second equation gives
$T_H=Q=0$.
\end{proof}

The use of $T_H$ in~\eqref{eq:homogeneous} denotes the retained
scaled coordinate; $T_{H,\mathrm{phys}}=y(H-P)/x(H-P)$ denotes
the physical ratio. No equation in this section divides by a variable
coordinate on the zero-input fiber.

\section{Three-coordinate numerator products and their ratios}
Fix a solution of~\eqref{eq:homogeneous} with $\beta\ne0$ and
$R\ne0$. Retain three field coordinates $k,F,Z$ with
\begin{equation}\label{eq:numerator}
                       F=k^2,\qquad Z^2+DF^2=R^2.
\end{equation}
For one further field coordinate $\Lambda$, impose
\begin{equation}\label{eq:ratio}
                              Z\Lambda=kS.
\end{equation}

\begin{proposition}\label{prop:numerator}
The point coordinates~\eqref{eq:numerator} describe points
$U\in2E(K)$, with the identity allowed, and satisfy
\begin{equation}\label{eq:producttest}
                         F=\beta^2 f(H)f(U).
\end{equation}
Every such point has the required lifts. Equation~\eqref{eq:ratio}
is exactly
\begin{equation}\label{eq:physicalratio}
                    \Lambda=\beta\frac{z_E(U)}{z_E(H)},
                    \qquad z_E=-x/y,
\end{equation}
with $z_E(O)=0$. On this $K$-locus both the reference denominator
and the numerator ordinate required in the formula are nonzero,
except that the numerator may be $O$, in which case $\Lambda=0$.
If $R=0$, the point equations instead force $F=Z=k=0$ and
\eqref{eq:ratio} does not constrain $\Lambda$.
\end{proposition}
\begin{proof}
The reference cover gives a root $r$ of $f(H)$; put
$r_\beta=\beta r$, so $r_\beta^2=R$. Define
\[
                 q_U=k/r_\beta,\qquad z_U=Z/R.
\]
Then
\[
 q_U^2=F/R,\qquad z_U^2+Dq_U^4=1.
\]
For $q_U\ne0$, the point
$U=(q_U^{-2},z_U/q_U^3)$ lies on $E$, has square abscissa,
and is in $2E(K)$ by Proposition~\ref{prop:image}.
For $q_U=0$, the two values $z_U=\pm1$ both represent $O$.
Conversely, a square root of $f(U)$ and its corresponding
$z_U=y(U)q_U^3$, including these two identity lifts, give the
displayed coordinates. Formula~\eqref{eq:producttest} follows from
$F=Rq_U^2$.

If $Z=0$, the point equations would give $F/R=\pm1/d$.
But $F/R=q_U^2$ is a square in $K$, contradicting its norm
$1/7$. Hence $Z\ne0$. Likewise a nonidentity reference point
on this cover has nonzero ordinate. Since $R\ne0$, $H\ne O$,
and its physical parameter is defined and nonzero. Using
$S=\beta^2rS_H$ and $z_E(U)=-q_U/z_U$ gives
\[
 \frac{kS}{Z}
 =\frac{\beta q_U(rS_H)}{r z_U}
 =\beta\frac{z_E(U)}{z_E(H)}.
\]
This remains zero at the identity numerator. Finally, if $R=0$,
the sum $Z^2+(dF)^2=0$ forces $Z=F=0$, and then $k=0$.
\end{proof}

The analogous construction for $H-P,V$ uses $Q,T_H$ and coordinates
$j,G,W,\Lambda_0$. On $RQ\ne0$ the two products used by the
detector are $F$ and $\overline G$. Their finite valuations at every
good odd place are even, directly from $F=k^2$ and $G=j^2$.
The comparison expressions
\[
             C=\Lambda-\overline{\Lambda_0},\qquad
             I=\Lambda-b_2
\]
are affine in retained coordinates. Thus their square coordinates
$C_2=C^2$, $I_2=I^2$ and the factor rows
\[
 wA_1=w^2+wC_2+F^2,\qquad
 wA_2=w^2+wI_2+\overline G^{\,2}
\]
are quadratic. The weighted products agree with those in
\cite{Quartic}; the two comparison expressions here are $\beta^2$
times its normalized comparisons. This observation does not transfer
the arithmetic proof of that construction: it does not encode the
difference $U-V$ or establish a complete defining system.

\section{Jacobi addition and the original trace}
The next presentation is a separate paired-point chart. It uses
unweighted reciprocal coordinates $f_U,f_V$.
Its trace is the function in the all-point height theorem
\cite[Theorem 1.1]{Trace}, and depends only on the difference $U-V$.

\begin{lemma}\label{lem:jacobi}
Let $q_U,z_U,q_V,z_V\in K$ satisfy
\[
 z_U^2+Dq_U^4=1,\qquad z_V^2+Dq_V^4=1.
\]
Let $U,V$ be their images on $E$, including the identity lifts.
Put $f_U=q_U^2$, $f_V=q_V^2$ and
\[
 N=q_Uz_V-q_Vz_U,\qquad C_0=1+Df_Uf_V.
\]
Then $C_0$ is nonzero, $N=0$ exactly when $U=V$, and for
$U\ne V$ one has
\begin{equation}\label{eq:jacobi}
                  x(U-V)=(C_0/N)^2.
\end{equation}
\end{lemma}
\begin{proof}
Both $f_U,f_V$ are squares, so $C_0>0$ at each real embedding.
If one $q$ is zero, its ordinate is $\pm1$, and $N=0$ forces
the other $q$ to be zero. Otherwise set $q_U=\lambda q_V$.
The equality $N=0$ gives $z_U=\lambda z_V$. The two curve equations
then give
\[
               (\lambda^2-1)(1+Dq_U^2q_V^2)=0.
\]
Thus $\lambda=\pm1$, and the simultaneous signs give the same
physical elliptic point. The converse follows from the same recovery
formulas.

For distinct abscissas, substituting
$x(U)=q_U^{-2}$, $y(U)=z_U/q_U^3$ and the corresponding values
for $V$ into the chord formula proves~\eqref{eq:jacobi}.
For clarity, putting $u=q_U,v=q_V,s=z_U,t=z_V$, the cleared
identity is
\begin{align*}
 &\bigl((sv^3+tu^3)^2-(u^2+v^2)(v^2-u^2)^2\bigr)(ut-vs)^2\\
 &\hspace{18mm}=u^2v^2(v^2-u^2)^2(1+Du^2v^2)^2,
\end{align*}
modulo $s^2+Du^4=t^2+Dv^4=1$.
If one point is $O$, the formula reduces to the other abscissa.
The remaining distinct-point equal-abscissa case is $V=-U$;
the numerator is then $\pm2q_Uz_U$, and the right side of
\eqref{eq:jacobi} is
$(1+Dq_U^4)^2/(4q_U^2z_U^2)$, exactly the duplication formula.
There is no nonidentity two-torsion point in the image over $K$,
so its denominator is nonzero. These cases exhaust the locus.
\end{proof}

Retain seven field coordinates
\[
                  h,f_U,f_V,z_U,z_V,Q_U,Q_V
\]
and two rational coordinates $T,P_h$. Write $h=h_0+\rho h_1$ and
impose
\begin{equation}\label{eq:tracecompression}
 T=2(h_0^2+2h_1^2),\qquad P_h=h_0h_1,
 \qquad H=T/2+2\rho P_h.
\end{equation}
Here $H$ is an affine field expression, not an additional coordinate.
Require the five field quadrics
\begin{equation}\label{eq:pairedjacobi}
 \begin{split}
 Q_U^2&=Hf_U,\qquad Q_V^2=Hf_V,\\
 z_U^2+Df_U^2&=1,\qquad z_V^2+Df_V^2=1,\\
 Q_Uz_V-Q_Vz_U&=1+Df_Uf_V.
 \end{split}
\end{equation}

\begin{theorem}\label{thm:paired}
On $h\ne0$, equations~\eqref{eq:tracecompression} and
\eqref{eq:pairedjacobi} describe ordered pairs
$U,V\in2E(K)$ with $U\ne V$, together with their square-root
lifts. They recover
\[
                     h^2=x(U-V),\qquad T=\Tr x(U-V).
\]
Every such ordered pair has a witness. The equations use sixteen
rational coordinates including $T$, and twelve rational quadrics.
The fiber $h=0$ consists exactly of
\begin{equation}\label{eq:zeroh}
 \begin{gathered}
 Q_U=Q_V=z_U=z_V=0,\qquad T=P_h=0,\\
 (f_U,f_V)=(1/d,-1/d)\quad\hbox{or}\quad(-1/d,1/d).
 \end{gathered}
\end{equation}
\end{theorem}
\begin{proof}
The rational rows give $H=h^2$ exactly. If $h\ne0$, set
$q_U=Q_U/h$ and $q_V=Q_V/h$. Then $f_U=q_U^2$ and
$f_V=q_V^2$, and the circle equations give the hypotheses of
Lemma~\ref{lem:jacobi}. The last equation is $hN=C_0$.
Its right side is nonzero, so $U\ne V$ and $h=C_0/N$.
The lemma gives the asserted abscissa and trace. Conversely choose
the two Jacobi lifts of any allowed pair and put $h=C_0/N$,
$Q_U=hq_U$, $Q_V=hq_V$, with $T,P_h$ defined by
\eqref{eq:tracecompression}. This supplies every row.

If $h=0$, then $H=0$ and $Q_U=Q_V=0$, so
$Df_Uf_V=-1$. The two circle equations give
\[
                    z_U^2+(df_Uz_V)^2=0.
\]
Formal reality forces both ordinates to vanish. Thus
$Df_U^2=Df_V^2=1$ and the signs are opposite, proving
\eqref{eq:zeroh}. Direct substitution verifies both solutions.
They correspond to the opposite nonzero-abscissa two-torsion
points on $E$; their difference is $T_0$, with abscissa zero.
The trace rows also force $T=P_h=0$.

Seven field coordinates contribute fourteen rational coordinates;
$T,P_h$ contribute two. Five field equations and two rational
equations give twelve rational quadrics.
\end{proof}

The zero-scale fiber is therefore present. Multiplying ratio equations
by $h$ can erase their information on that fiber and requires a separate
check in an application. If an application independently forces
$f_U,f_V$ to be squares in $K$, both points in~\eqref{eq:zeroh}
are excluded by the norm obstruction above.

The point image also fixes the appropriate integer label. Relative to
$P=2R_0$, the difference in Theorem~\ref{thm:paired} has an integer
index. Thus an application to all integer inputs may take
$\beta=X$ and $A=[X]P$. The assignment $\beta=X/2$ and
$A=[X]R_0$ would omit odd inputs from this particular paired cover.

\section{Coordinate costs and the remaining interface}
The proven presentations have the following separate counts.
\begin{center}
\begin{tabular}{lrr}
\toprule
Presentation & Rational coordinates & Rational quadrics\\
\midrule
Three-coordinate reference cover &6&4\\
Four-coordinate reference cover &8&6\\
One weighted numerator and its ratio &8&6\\
Jacobi pair with original rational trace &16&12\\
\bottomrule
\end{tabular}
\end{center}
The homogeneous reference rows use the same eight reference
coordinates and the shared input powers $b_2,b_4$; their definition
must be counted wherever those powers are not already retained.
The last row includes the original trace and the one additional
rational product coordinate $P_h$.

These counts are not a census of a combined defining system.
The weighted numerator chart uses $F=Rf_U$, $G=Qf_V$ and
$Z=Rz_U$, $W=Qz_V$, whereas Theorem~\ref{thm:paired} uses
the unweighted point coordinates. With $k^2=F$, $j^2=G$, the
Jacobi addition interface becomes
\[
 h\bigl(\sqrt R\,kW-\sqrt Q\,jZ\bigr)=RQ+DFG.
\]
The square roots exist on the nonzero reference locus, but this
display is not a quadratic equation in the retained weighted
coordinates. Its products, the original difference trace and its
boundary conditions must all be converted before a full variable
saving can be asserted. No theorem here supplies that conversion.

\section{Exact verification}\label{sec:verification}
The package
\nolinkurl{computation/elliptic_charts/reference_covers_v47}
contains the universal exact checker, its certificate, and a separate
independent checker. The universal checker verifies the two reference
maps, both compositions, the compact-reference identities, all four
boundary points, homogeneous input scaling, the numerator product and
ratio identities, the Jacobi difference identity, trace compression,
and the zero-height sum-of-squares identity. It also authenticates
the specialized quadratic-field constants and the stated block counts.

The independent checker reconstructs the reference maps and performs
its own Buchberger calculation with provenance in the original
defining ideals. It checks that every certificate denominator is
supported on $e,c,e-c,e+c$. Its scope is the reference isomorphism
and boundary statements; it does not independently certify a complete
defining polynomial. Both programs print their results by default and
write a receipt only when an explicit output path is supplied.

\begin{thebibliography}{9}
\bibitem{Three}
\emph{A three-coordinate quadratic chart for an elliptic point and its translate},
\nolinkurl{papers/elliptic_charts/three_coordinate_reference_20261115_v40.tex}.
\bibitem{Compact}
\emph{Quadratic rescaling of elliptic reference and height coordinates},
\nolinkurl{papers/elliptic_charts/scaled_reference_height_20261116_v43.tex}.
\bibitem{Trace}
\emph{Trace heights on quadratic twists with $j=1728$},
\nolinkurl{papers/trace_heights/trace_on_all_points_20261113_v25.tex}.
\bibitem{Quartic}
\emph{A quartic in 74 variables from weighted point products},
\nolinkurl{papers/quartic_existence/quartic_abc74_20261116_v46.tex}.
\end{thebibliography}
\end{document}
